\chapter{SPICE-Netzlisten}
\label{kap:spice}

esim liest SPICE auf zwei Wegen:

\begin{itemize}
  \item \textbf{Modellkarten und Teilschaltungen} direkt in der
    Eingabedatei (Kapitel~\ref{kap:format-modelle} und
    Abschnitt~\ref{sec:spice-subckt}). So werden Modelle von Herstellern
    unverändert benutzt.
  \item \textbf{Ganze Netzlisten} übersetzt das Werkzeug
    \datei{tools/spice2json.py} einmalig in das JSON-Format. Danach lassen
    sie sich im Editor öffnen, zeichnen und weiterbearbeiten.
\end{itemize}

\section{spice2json}

\begin{shell}
python3 tools/spice2json.py examples/spice/amplifier.cir -o amplifier.json
esim amplifier.json
\end{shell}

Ohne \code{-o} schreibt das Werkzeug auf die Standardausgabe. Ein Beispiel
für eine Netzliste:

\begin{spicecode}
Common-emitter amplifier with an RC filter subcircuit (SPICE netlist)
.param rload=10k cfilt=1n

VCC vcc 0 DC 12
Vin in 0 DC 0 AC 1 SIN(0 10m 1k)
Cin in b 10u
R1 vcc b 47k
R2 b 0 10k
Q1 c b e QGEN
RC vcc c 4.7k
RE e 0 1k
CE e 0 100u
Cout c out 10u
RL out 0 {rload}

X1 out filt rcfilter c={cfilt}
D1 filt clamp DGEN
Vclamp clamp 0 DC 3

.subckt rcfilter a b r=1k c=10n
Rf a b {r}
Cf b 0 {c}
.ends

.model QGEN NPN(Is=1e-14 Bf=200 Br=2 Vaf=100
+ Ikf=0.3 Rb=10 Rc=1 Cje=10p Cjc=5p
+ Tf=0.5n Tr=50n)
.model DGEN D(Is=2n Rs=0.5 N=1.75 Cjo=4p M=0.4 Tt=20n)

.op
.ac dec 20 10 10meg
.noise v(out) Vin dec 10 10 100k
.tran 10u 5m
.four 1k v(out)
.print tran v(in) v(out) v(filt)
.end
\end{spicecode}

\section{Was übersetzt wird}

\begin{center}
\begin{tabularx}{\textwidth}{@{}lL@{}}
  \toprule
  Art & Unterstützt \\
  \midrule
  Bauteile & \code{R C L K V I E G F H D Q M J X S T}, \code{O} (verlustbehaftete Leitung mit LTRA-Modell) \\
  Quellenformen & \code{DC}, \code{AC}, \code{SIN}, \code{PULSE}, \code{PWL} \\
  Anweisungen & \code{.model}, \code{.subckt}/\code{.ends}, \code{.param}, \code{.op}, \code{.dc}, \code{.ac}, \code{.tran}, \code{.noise}, \code{.four}, \code{.ic}, \code{.nodeset}, \code{.temp}, \code{.options}, \code{.print} \\
  Werte & SPICE-Zahlen (\code{10k}, \code{4.7meg}, \code{1m} = Milli), Ausdrücke in \code{{...}} als Parameter \\
  \bottomrule
\end{tabularx}
\end{center}

\begin{itemize}
  \item Namen und Knoten werden kleingeschrieben, weil SPICE nicht zwischen
    Groß- und Kleinschreibung unterscheidet. Ausgaben heißen danach etwa
    \code{v(out)} und \code{i(vin)}.
  \item Modelle bleiben als SPICE-Zeilen erhalten; esim liest sie selbst.
  \item \code{.print} wird zu \code{"outputs"}, \code{.four} zum Block
    \code{fourier} der Transientenanalyse, ein Druckschritt von
    \code{.tran} zu \code{output_step}.
  \item Was nicht übersetzt werden kann (etwa \code{.tf} oder unbekannte
    Bauteile), wird gemeldet und ausgelassen. Eingebundene Dateien
    (\code{.include}, \code{.lib}) werden nicht gelesen; ihren Inhalt
    einfügen oder die Modelle über die Bibliothek des Editors einbinden.
\end{itemize}

\section{Unterschiede zu SPICE}

\begin{center}
\begin{tabularx}{\textwidth}{@{}P{4.2cm}L@{}}
  \toprule
  Thema & in esim \\
  \midrule
  Präfixe im JSON & Groß/klein zählt: \code{M} Mega, \code{m} Milli (in SPICE-Text wie SPICE). \\
  Zweigströme & \code{i(V1)} wie SPICE: vom Plus-Knoten durch die Quelle. \\
  Integration & Standard ist BDF variabler Ordnung, nicht Trapez; \code{"method": "trap"} wählt die Trapezregel. \\
  Schrittweite & \code{tstep} begrenzt die Schrittweite; der Druckschritt von \code{.tran} wird zu \code{output_step}. \\
  Schalter & Ideal, mit exakt bestimmtem Umschaltzeitpunkt; PSpice-\code{VSWITCH} als stetiger Übergang. \\
  Ableitungen & Exakt durch automatisches Differenzieren, auch in Verhaltensquellen. \\
  Modelle & Diode, Gummel-Poon, JFET, MOS 1/3, BSIM3v3; dazu EKV. BSIM4 und neuere fehlen. \\
  \bottomrule
\end{tabularx}
\end{center}

\section{Prüfen}

\datei{tools/test_spice2json.py} übersetzt
\datei{examples/spice/amplifier.cir} und vergleicht das Ergebnis mit einer
von Hand geschriebenen Netzliste derselben Schaltung. Der Test läuft mit
\code{ctest} (Gruppe \code{spice_import}).
